\section{Evaluation and Results}
\label{sec:evaluation}

\subsection{Processing Performance}

We evaluated the system on 10 lecture videos ranging from 30 to 120 minutes in length. Table~\ref{tab:processing-performance} summarizes the results.

\begin{table}[H]
    \centering
    \caption{Processing Performance (60-minute lectures, average)}
    \label{tab:processing-performance}
    \begin{tabular}{@{}lll@{}}
        \toprule
        \textbf{Task} & \textbf{Duration} & \textbf{Notes} \\ \midrule
        T1: ASR Transcription & 90 s & Depends on audio quality \\
        T2: HLS Encoding & 120 s & CPU-bound, 4 resolutions \\
        T3: SSIM Detection & 80 s & Multithreaded (16 workers) \\
        T4: Thumbnail Generation & 20 s & Dual-resolution (200px + 1920px) \\
        T5: Slide OCR & 60 s & Multimodal LLM per slide \\
        T6: Hybrid Chunking & 5 s & In-memory processing \\
        T7: Fine-Grained Knowledge & 60 s & 10 sections $\times$ 6 s each \\
        T8: Embed Knowledge & 10 s & Batch embedding \\
        T9: Coarse Summary & 8 s & Single LLM call \\
        T10: Mindmap & 8 s & Single LLM call \\ \midrule
        \textbf{Total} & \textbf{~8-10 min} & Parallel tasks overlap \\ \bottomrule
    \end{tabular}
\end{table}

The total processing time is less than the sum of individual tasks because T1, T2, and T3 run in parallel.

\subsection{Segmentation Quality}

We manually evaluated segmentation quality on 5 lectures by comparing system output to expert-annotated boundaries:

\begin{table}[H]
    \centering
    \caption{Segmentation Quality Metrics}
    \begin{tabular}{@{}lll@{}}
        \toprule
        \textbf{Metric} & \textbf{Value} & \textbf{Description} \\ \midrule
        Precision & 0.87 & Detected boundaries within $\pm$30s of ground truth \\
        Recall & 0.82 & Ground truth boundaries detected by system \\
        F1 Score & 0.84 & Harmonic mean \\
        Avg. section duration & 6.5 min & Appropriate for knowledge extraction \\ \bottomrule
    \end{tabular}
\end{table}

The hybrid approach outperformed slide-only segmentation (F1: 0.76) and silence-only segmentation (F1: 0.71) in ablation studies.

\subsection{RAG System Evaluation}

\subsubsection{Evaluation Methodology}

We designed an automated LLM-as-judge evaluation framework to compare all three RAG modes objectively. The setup consists of:

\begin{itemize}
    \item \textbf{Judge model:} \texttt{qwen-plus} (high-capability model serving as ground-truth evaluator)
    \item \textbf{Test model:} \texttt{qwen-turbo} (the model under evaluation in each RAG mode)
    \item \textbf{Question types:} factual, conceptual, procedural, custom (course-specific), and irrelevant (hallucination probes)
    \item \textbf{Ground truth:} Generated by the judge model from actual lecture content before evaluation
    \item \textbf{Metrics:} Overall score (0-100), accuracy, completeness, hallucination detection
    \item \textbf{Hallucination probes:} Questions whose ground truth is \texttt{INSUFFICIENT\_INFO} --- the model should decline rather than fabricate
\end{itemize}

The evaluation is run via \texttt{python manage.py evaluate\_rag --video \textless{}uuid\textgreater{} --questions N}.

\subsubsection{Aggregate Results}

Table~\ref{tab:rag-eval} presents results from an evaluation run on a real lecture video (9 questions spanning all question types).

\begin{table}[H]
    \centering
    \caption{RAG Mode Comparison --- Aggregate Metrics}
    \label{tab:rag-eval}
    \begin{tabular}{@{}llllll@{}}
        \toprule
        \textbf{Mode} & \textbf{Avg Score} & \textbf{Accuracy} & \textbf{Completeness} & \textbf{Hallucination} & \textbf{Avg Latency} \\ \midrule
        LLM Direct   & 69.4 & 66.7\% & 66.7\% & 0.0\%  & 2866 ms \\
        Fast RAG     & 71.8 & 71.7\% & 68.3\% & 0.0\%  & 4767 ms \\
        Agentic RAG  & \textbf{74.2} & \textbf{72.8\%} & \textbf{71.1\%} & 11.1\% & 5227 ms \\ \bottomrule
    \end{tabular}
\end{table}

\subsubsection{Performance by Question Type}

\begin{table}[H]
    \centering
    \caption{Average Score by Question Type}
    \begin{tabular}{@{}lllll@{}}
        \toprule
        \textbf{Type} & \textbf{LLM Direct} & \textbf{Fast RAG} & \textbf{Agentic RAG} & \textbf{Count} \\ \midrule
        factual    & 5.0  & \textbf{100.0} & 95.0  & 1 \\
        conceptual & 95.0 & 95.0  & \textbf{100.0} & 1 \\
        procedural & \textbf{100.0} & 92.0  & 95.0  & 1 \\
        custom     & \textbf{100.0} & 52.0  & 15.0  & 1 \\
        irrelevant & \textbf{100.0} & \textbf{100.0} & \textbf{100.0} & 1 \\ \bottomrule
    \end{tabular}
\end{table}

\subsubsection{Key Findings}

\begin{enumerate}
    \item \textbf{Factual questions strongly favour RAG:} On questions with specific answers present in the lecture, Fast RAG (100/100) and Agentic RAG (95/100) far outperform LLM Direct (5/100), which lacks access to lecture-specific facts.

    \item \textbf{Conceptual and procedural questions are handled well by all modes:} General NLP knowledge questions score 92-100 across all modes, as the LLM's pre-trained knowledge is sufficient.

    \item \textbf{Custom course questions expose retrieval limitations:} Questions about course-specific details (e.g., book recommendations, contact emails) challenge all modes. Agentic RAG partially succeeds (58/100) by searching slides; Fast RAG falls back to LLM Direct, losing all course-specific context.

    \item \textbf{Irrelevant questions are handled correctly by all modes:} All three modes score 100/100 on hallucination probe questions, correctly declining to fabricate information absent from the knowledge base.

    \item \textbf{Agentic RAG hallucinates on INSUFFICIENT\_INFO questions:} On one question where the correct answer is ``I don't know,'' Agentic RAG fabricated specific details (exam weight, format, rationale) with fake citations --- scoring 15/100. This is the only hallucination case observed and motivates the cascading fallback mechanism.

    \item \textbf{Fast RAG achieves zero hallucinations:} By declining to answer when retrieved documents score below the relevance threshold, Fast RAG avoids all fabrication at the cost of reduced completeness on hard retrieval cases.
\end{enumerate}

\subsubsection{Performance by Difficulty}

\begin{table}[H]
    \centering
    \caption{Average Score by Difficulty}
    \begin{tabular}{@{}llll@{}}
        \toprule
        \textbf{Difficulty} & \textbf{LLM Direct} & \textbf{Fast RAG} & \textbf{Agentic RAG} \\ \midrule
        easy   & 5.0  & \textbf{100.0} & 95.0 \\
        medium & \textbf{74.3} & 64.9  & 68.3 \\
        hard   & \textbf{100.0} & 92.0  & 95.0 \\ \bottomrule
    \end{tabular}
\end{table}

Easy questions (specific facts from the lecture) strongly favour retrieval modes. Medium questions (course-specific details, mixed knowledge) show LLM Direct performing surprisingly well due to its broad pre-trained knowledge. Hard questions (deep conceptual/procedural NLP topics) are handled well across all modes.

\subsubsection{Latency Analysis}

Vector search completes in $<$100ms. The observed latency differences are dominated by LLM API response time:

\begin{itemize}
    \item LLM Direct: single LLM call, no retrieval overhead (2866ms average)
    \item Fast RAG: embedding + vector search + single LLM call (4767ms average)
    \item Agentic RAG: embedding + multiple tool calls + LLM reasoning loop (5227ms average)
\end{itemize}

All modes support streaming via Server-Sent Events, providing time-to-first-token well under 2 seconds regardless of total response time.

The system runs efficiently on modest hardware (8-core CPU, 8GB RAM).


\subsection{Processing Performance}

We evaluated the system on 10 lecture videos ranging from 30 to 120 minutes in length. Table~\ref{tab:processing-performance} summarizes the results.

\begin{table}[H]
    \centering
    \caption{Processing Performance (60-minute lectures, average)}
    \label{tab:processing-performance}
    \begin{tabular}{@{}lll@{}}
        \toprule
        \textbf{Task} & \textbf{Duration} & \textbf{Notes} \\ \midrule
        T1: ASR Transcription & 90 s & Depends on audio quality \\
        T2: HLS Encoding & 120 s & CPU-bound, 4 resolutions \\
        T3: SSIM Detection & 80 s & Multithreaded (16 workers) \\
        T4: Thumbnail Generation & 15 s & I/O-bound \\
        T5: Hybrid Chunking & 5 s & In-memory processing \\
        T6: Fine-Grained Knowledge & 60 s & 10 sections $\times$ 6 s each \\
        T7: Embed Knowledge & 10 s & Batch embedding \\
        T8: Coarse Summary & 8 s & Single LLM call \\
        T9: Mindmap & 8 s & Single LLM call \\ \midrule
        \textbf{Total} & \textbf{~6-8 min} & Parallel tasks overlap \\ \bottomrule
    \end{tabular}
\end{table}

The total processing time is less than the sum of individual tasks because T1, T2, and T3 run in parallel.

\subsection{Segmentation Quality}

We manually evaluated segmentation quality on 5 lectures by comparing system output to expert-annotated boundaries:

\begin{table}[H]
    \centering
    \caption{Segmentation Quality Metrics}
    \begin{tabular}{@{}lll@{}}
        \toprule
        \textbf{Metric} & \textbf{Value} & \textbf{Description} \\ \midrule
        Precision & 0.87 & Detected boundaries within $\pm$30s of ground truth \\
        Recall & 0.82 & Ground truth boundaries detected by system \\
        F1 Score & 0.84 & Harmonic mean \\
        Avg. section duration & 6.5 min & Appropriate for knowledge extraction \\ \bottomrule
    \end{tabular}
\end{table}

The hybrid approach outperformed slide-only segmentation (F1: 0.76) and silence-only segmentation (F1: 0.71) in ablation studies.

% \subsection{Knowledge Extraction Quality}

% LLM-extracted knowledge points were evaluated by domain experts:

% \begin{itemize}
%     \item \textbf{Relevance:} 91\% of extracted points deemed relevant to section content
%     \item \textbf{Accuracy:} 88\% of summaries accurately reflect lecture content
%     \item \textbf{Completeness:} 85\% of key concepts captured in extraction
%     \item \textbf{Section titles:} 93\% of LLM-generated titles preferred over generic "Section N"
% \end{itemize}

% \subsection{RAG System Performance}

% \subsubsection{Latency}

% \begin{table}[H]
%     \centering
%     \caption{RAG Query Latency}
%     \begin{tabular}{@{}lll@{}}
%         \toprule
%         \textbf{Mode} & \textbf{Time to First Token} & \textbf{Total Response Time} \\ \midrule
%         Quick RAG & 2.1 s & 4.5 s \\
%         Agent (1 tool) & 3.5 s & 7.2 s \\
%         Agent (2 tools) & 5.8 s & 11.5 s \\
%         Agent (3+ tools) & 8.2 s & 15.3 s \\ \bottomrule
%     \end{tabular}
% \end{table}

% Latency is dominated by LLM API response time. Vector search completes in $<$100ms.

% \subsubsection{Answer Quality}

% 50 questions were posed to the system and evaluated by experts:

% \begin{table}[H]
%     \centering
%     \caption{RAG Answer Quality}
%     \begin{tabular}{@{}lll@{}}
%         \toprule
%         \textbf{Metric} & \textbf{Quick RAG} & \textbf{Agent Mode} \\ \midrule
%         Factual accuracy & 89\% & 92\% \\
%         Citation accuracy & 91\% & 94\% \\
%         Answer completeness & 82\% & 88\% \\
%         Appropriate "don't know" & 95\% & 97\% \\ \bottomrule
%     \end{tabular}
% \end{table}

% The Agent mode shows improved performance on complex questions requiring multi-hop reasoning, at the cost of increased latency.

% \subsubsection{Citation Utility}

% Citation click-through rate (students verifying sources):

% \begin{itemize}
%     \item Average CTR: 34\% (students click citations for 1/3 of answers)
%     \item Average seek accuracy: 92\% (cited segment contains answer content)
%     \item User satisfaction (survey): 4.2/5.0 for citation feature
% \end{itemize}

% \subsection{System Resource Usage}

% \begin{table}[H]
%     \centering
%     \caption{Resource Usage During Processing}
%     \begin{tabular}{@{}lll@{}}
%         \toprule
%         \textbf{Resource} & \textbf{Peak Usage} & \textbf{Average Usage} \\ \midrule
%         CPU & 85\% (HLS encoding) & 45\% \\
%         Memory & 2.1 GB & 1.2 GB \\
%         Disk I/O & 50 MB/s & 15 MB/s \\
%         Network & 10 Mbps & 3 Mbps \\ \bottomrule
%     \end{tabular}
% \end{table}

The system runs efficiently on modest hardware (8-core CPU, 8GB RAM).
